\section{The Worship of Heaven}

In John's vision the door is opened, the throne is revealed, and the life surrounding it is gathered into adoration (Rev 4). The imagery is rich with light, living creatures, elders, vesture, and crowns. Its center is the one seated upon the throne. The living creatures acclaim God's holiness, and the elders fall before him. We can allow its many figures to awaken wonder without assigning each a physical arrangement or claiming to know every detail of the life they signify. What is clear is the joyful acknowledgment of God's holiness and of the dependence of every created thing upon his will.

The elders cast their crowns before the throne. A crown ordinarily signifies distinction, authority, or a victory already attained; here it becomes an offering. They honor God with the crowns they have received. They acknowledge that their glory comes from him and offer it in thanksgiving. An earthly achievement can teach us a little of this movement when success awakens thankfulness for the people, opportunities, and powers through which it became possible. Remembering how much we have received can correct the wish to regard an achievement as wholly our own.

The elders praise God because all things exist through his will. Before we consider what a creature has done for us, it already has its place within his generosity. This can enlarge our prayer beyond the circle of immediate benefit. We can give thanks for a truth we may never master, a beauty far from our own country, or a person whose service will never become known to us. The world contains goods that are real without becoming our possessions. To worship the Creator is to rejoice that he gives being so abundantly. Such thanksgiving can continue even when our own needs remain pressing, because God's goodness extends beyond what we have received ourselves.

John's vision then concentrates upon a sealed book and the grief that no one can open it (Rev 5). He is told of the victory of the lion of the tribe of Judah, and he sees a Lamb standing as though slain. The worship of heaven is inseparable from the saving death of Christ. The Lamb receives the book, and the new canticle praises him for redeeming a people by his blood. His self-offering remains the subject of their praise in glory.

We can desire freedom from suffering while resisting the costly love through which Christ draws us to himself. The promise is indeed freedom from suffering, but the life that reaches fulfillment is the life given to us through Christ's self-offering. We shall rejoice in a love whose earthly form included obedience, patience, and the gift of himself for those he came to save. We cannot prepare for that joy by treating every costly demand of charity as foreign to happiness. In offering something costly for another's good, we seek to share the charity for which the Lamb is praised.

The new canticle praises the Lamb for redeeming a people by his blood from every tribe, tongue, people, and nation. The breadth of the communion is part of its glory. Each of us enters the Church within particular languages, customs, and relationships; none of these beginnings encloses the whole company of the redeemed. Heavenly hope teaches us to receive that company with gratitude. The faith of Christians formed in unfamiliar circumstances is a gift to the same body. Their sufferings matter to us, and their holiness belongs to a communion we desire to share. Concern confined to Christians whose circumstances resemble our own falls short of the company the canticle describes.

The redeemed are described as a kingdom and priests. Worship is thus bound to the dignity Christ gives his people. They participate in the praise offered through Christ and receive their dignity from him. Their life under his sovereignty exceeds what present political arrangements can describe. Whatever differences of service and glory belong to the saints, all are joined in acknowledging him. Each contributes to their common praise as someone whom Christ has redeemed.

John sees bowls of incense identified with the prayers of the saints. Within the vision's splendor there is room for prayer: the cry that rises to God is received within the worship of heaven. This image can sustain the Christian who has little outward evidence that his intercession matters. A prayer offered in a crowded household, from a prison, or during an ordinary journey may be hidden from every observer. Its value is known to God. John identifies the incense with their prayers without describing how individual saints know particular requests. The petition hidden on earth is received by the God to whom the saints offer praise.

The praise widens through myriads of angels and finally through every creature in heaven, on earth, and in the sea. Glory is given to the one on the throne and to the Lamb. We shall be fulfilled within a good in which others also rejoice. Their worship will belong to our happiness because the God they praise is the God we love. Even now we can learn something of this when another person's thanksgiving enlarges our own, or when the Church's common prayer gives words to a gratitude we could not have expressed alone.

The contemplation of heavenly worship prepares us for the worship entrusted to us now. It can teach us to enter the Church's praise with attention, to receive silence as room for adoration, and to offer our voices as members of a body. The Father is worthy of love before we have obtained an answer to each particular petition. Christ is worthy of thanksgiving because he has redeemed us. We can adore while a petition remains unanswered, joining the praise offered for God's holiness and Christ's accomplished redemption. Our worship in faith joins that acknowledgment and awakens the desire to take our place, wholly and forever, in its fulfillment.

